% !TeX program = xelatex
\documentclass{cumcmthesis}
% \documentclass[withoutpreface,bwprint]{cumcmthesis} %去掉封面与编号页，电子版提交的时候使用。


\usepackage[framemethod=TikZ]{mdframed}
\usepackage{url}   % 网页链接
\usepackage{subcaption} % 子标题、子图、子表格
\usepackage{pifont}%提供大量的特殊符号、图标
\usepackage{listings}%代码高亮包
\usepackage{float} %管控图片、图表浮动体
%% 下面是参赛的基本信息填写
\title{全国大学生数学建模竞赛编写的 \LaTeX{} 模板}
\tihao{A}
\baominghao{xxxx}
\schoolname{XX大学}
\membera{ }
\memberb{ }
\memberc{ }
\supervisor{ }
\yearinput{2026}
\monthinput{09}
\dayinput{13}

\begin{document}

 \maketitle  % 输出标题区（标题 / 作者 / 日期），必须在 document 环境内
 \begin{abstract}
随着xx的发展,xx问题是xx中的重要研究课题（或xx现象日益严重）。本文针对xx中的问题，基于xx和xx思想,通过确定xx、xx、xx等指标，以xx、xx为目标建立了xx模型，并使用xx算法对模型进行求解。
%模板中只是提醒需要注意，而不是真的需要加粗字体
\textbf{针对问题一，}

\textbf{针对问题二，}

\textbf{针对问题三，}


最后,我们对提出的模型进行全面的评价：本文的模型贴合实际,能合理解决提出的问题,具有实用性强,算法效率高等特点,该模型在xx,xx,xx方面也能使用。

\textbf{关键词： }关键词1  关键词2  关键词3  关键词4  关键词5
%此处是提醒同学可以查看具体的网址链接，而不是真的需要贴一个链接
\href{https://www.mcm.edu.cn/html_cn/node/4cd596519c9eb9fbd866398f6df0caa3.html}{全国大学生数学建模竞赛论文格式规范（2026年修订稿）}

\end{abstract}


\begin{enumerate} 
\item 标题：写成“基于××模型/方法的××问题研究”，也可根据论文内容自行拟定。
\item 摘要结构：建议采用“背景—方法—结果—结论”四步法。
\item 篇幅要求：摘要与关键词原则上控制在一页内，不写英文摘要。
\item 关键词：一般选取3～5个，以核心模型、关键算法和研究对象为主。
\item 写作原则：按照题目各问依次展开，重点突出模型、算法和具体结果。
\end{enumerate}

论文摘要不是论文目录，也不是正文的简单压缩。

第一步：背景

简要说明研究对象、问题背景和主要任务，通常用1～2句话完成。背景不宜过长，不要大段照抄题目。

参考句式：针对××背景下的××问题，本文围绕××、××和××等任务展开研究。

第二步：方法

按照问题一、问题二、问题三的顺序，分别说明所采用的数据处理方法、数学模型和求解算法。

参考句式：针对问题一，首先对××数据进行预处理，选取××指标，\textbf{建立××模型}，并\textbf{采用××算法}进行求解。
注意写出具体的模型或算法名称，避免只写“建立数学模型”“利用软件求解”。

第三步：结果

给出各问题的关键答案，包括主要数值、排序、预测结果、最优方案、误差或评价指标。结果应与前面的方法一一对应。

参考句式：结果表明，\textbf{最优方案为××}，\textbf{其目标函数值为××}；\textbf{模型预测误差为××\%}。
避免使用“结果较好”“模型合理”“得到满意结果”等模糊表述。

第四步：结论

简要说明模型的准确性、稳定性、优缺点及推广价值，最好结合误差分析、灵敏度分析或对比实验进行评价。

参考句式：灵敏度分析表明，主要参数在一定范围内变化时，模型结果保持稳定，可进一步推广至××问题。

每个问题尽量按照“问题任务＋模型方法＋求解结果”组织，整体做到：

\textbf{背景要短，方法要准，结果要实，结论要有依据。}

官方要求摘要页包含论文题目、摘要和关键词，原则上不超过一页。



\newpage

%% ------------------------------------------------------------
%% 正文各节（Writer 分节生成，顺序固定，不得增删调序）
%% 以下各节内容拆分自 example.tex（原版本），源文件 example.tex 保持不动
%% ------------------------------------------------------------
\section{问题重述}
\subsection{问题背景}
\begin{enumerate}[label=(\arabic*), leftmargin=2.5em, itemsep=0.3em]
\item \textbf{不要直接复制赛题。} 应提炼题目的背景、已知条件和研究任务，用自己的语言重新组织表述。直接复制原题容易导致论文重复率较高。

\item \textbf{删除冗余内容。} 题目中的故事背景、案例介绍、新闻材料等一般无需全部保留，只保留与建模相关的信息即可。

\item \textbf{保持原题意思不变。} 不要随意增加题目没有给出的条件，也不要提前加入模型假设或个人理解。

\item \textbf{按题目顺序整理。} 将问题一、问题二、问题三依次列出，使后续"问题分析"和"模型建立"保持一致。

\item \textbf{控制篇幅。} 一般控制在\textbf{半页至一页}即可，不必花费过多篇幅。
\end{enumerate}

\textbf{问题重述的目的，是让读者快速了解题目，而不是展示建模能力，因此做到"准确、简洁、完整"即可。}

\subsection{问题提出}

%\paragraph{问题一：} 不会缩进 

\textbf{问题一：}

\textbf{问题二：}

\textbf{问题三：}
   % 一、问题重述
\section{问题分析}
\subsection*{问题一：}
%也可以考虑使用\paragraph{问题一：} 不占用章节编号层级
\begin{enumerate}[label=(\arabic*), leftmargin=2.5em, itemsep=0.3em]
\item \textbf{问题分析是论文的重要部分。} 主要分析每个问题的特点、难点及建模思路，为后续模型建立提供依据。

\item \textbf{一般按问题分别分析。} 建议按照问题一、问题二、问题三依次展开，每一问对应一份分析。

\item \textbf{重点分析思路。} 说明为什么采用该模型、哪些因素需要重点考虑、模型如何解决问题，而不是直接推导公式。

\item \textbf{建立问题之间的联系。} 若后续问题依赖前一问的结果，应说明各问题之间的数据、参数或模型关系。

\item \textbf{不写求解结果。} 本部分只分析建模思路，不提前给出计算结果、实验结果或最终结论。
\end{enumerate}

\textbf{问题分析的目的，是告诉评阅老师："为什么这样建模"，而不是"模型怎么算"。}

\textbf{问题分析主要写什么？}

对于每一问，建议围绕以下几个方面展开：
\begin{enumerate}[label=(\arabic*), leftmargin=2.5em, itemsep=0.3em]
\item \textbf{明确问题类型} 判断属于评价、预测、分类、优化、决策、机理分析等哪一类问题。

\item \textbf{分析关键因素} 找出影响结果的主要变量、约束条件和题目提供的数据，说明哪些因素需要重点考虑。

\item \textbf{确定建模思路} 结合问题特点，分析适合采用哪些数学模型或算法，并简要说明选择理由。

\item \textbf{说明求解目标} 明确模型最终需要输出什么结果，如最优方案、预测值、评价结果、分类结果等。
\end{enumerate}

\textbf{写作提示}

\begin{enumerate}[label=(\arabic*), leftmargin=2.5em, itemsep=0.3em]
\item 先分析问题，再选择模型，不要为了使用某个模型而分析问题。

\item 模型名称不宜罗列过多，应选择最符合题意的方法，并说明选择依据。

\item 每一问尽量围绕"问题特点—关键因素—建模思路—预期结果"展开，保持逻辑清晰。

\item 问题分析可以在完成模型建立后再回过头完善，但内容必须与正文保持一致。
\end{enumerate}

\textbf{写作原则：分析问题、明确思路、说明方法、不写结果。}

\subsection*{问题二：}

\subsection*{问题三：}
      % 二、问题分析
\section{模型假设}
\subsection{模型基本假设}
\begin{enumerate}[label=(\arabic*), leftmargin=2.5em, itemsep=0.3em]
\item \textbf{假设应服务于建模。} 模型假设用于简化实际问题、明确研究范围，并为模型建立提供合理前提。

\item \textbf{假设必须有依据。 }可以忽略影响较小的因素，也可以将短期内变化不明显的量视为常数，但应符合实际情况。

\item \textbf{只写必要假设。} 不要罗列与模型无关、对求解没有影响的内容，也不要为了凑篇幅设置过多假设。

\item \textbf{假设应具体明确。} 避免使用“其他条件均不变”“不考虑特殊情况”等过于笼统的表述，应说明具体忽略了什么、保持什么不变。

\item \textbf{后文保持一致。} 模型建立、参数设置和结果分析必须遵循已提出的假设；若关键假设可能影响结果，应进行误差分析或灵敏度分析。
\end{enumerate}
\textbf{模型假设不是随意简化，而是在保证主要规律不变的前提下，忽略次要因素。}

\textbf{常见假设类型}

\ding{172} \textbf{忽略次要因素}
当某些因素对研究结果影响较小时，可以暂不考虑。
假设在研究时间范围内，天气变化对运输时间的影响可以忽略。

\ding{173} \textbf{将变量视为常数}
对于短期内变化较小或难以获取的数据，可在一定范围内视为常数。
假设研究期间产品单价和运输成本保持不变。

\ding{174} \textbf{理想化研究条件}
对实际过程进行合理理想化，以降低模型复杂度。
假设设备运行状态正常，不发生突发故障。

\ding{175} \textbf{排除小概率事件}
在题目未特别要求时，可不考虑极端情况或偶发事件。
假设研究期间不发生重大自然灾害和突发公共事件。

\textbf{写作提示}
\begin{enumerate}[label=(\arabic*), leftmargin=2.5em, itemsep=0.3em]
\item 每条假设尽量只表达一个意思，并采用编号分条书写。

\item 不要把题目已经明确给出的条件重复写成假设。

\item 避免使用明显不符合实际、只是为了方便计算的强假设。
\end{enumerate}


\textbf{写作原则：合理、必要、具体、可解释。}
   % 三、模型假设
\section{符号说明}
\begin{table}[!htbp]%[!htbp] 是浮动体位置优先级。! 表示“强制忽略LaTeX内部浮动静默限制”，h（here，这里）、t（top，页面顶部）、b（bottom，页面底部）、p（单独成页）。
    \caption{符号说明}\label{tab:demo} \centering %\caption 生成表头标题，\label 生成引用标签，\centering 让表格在页面水平居中。
    \renewcommand{\arraystretch}{1.0}%默认行高的一倍，可以依据公式情况调整
    \begin{tabularx}{0.7\textwidth}{%表格总宽度占页面宽度的70%（给右侧留白批注空间）。
        >{\centering\arraybackslash}X %三列均为“自动拉伸且居中”的弹性列。
        >{\centering\arraybackslash}X
        >{\centering\arraybackslash}X
    }
        \toprule[1.5pt]%三线表中顶部最粗，底部次之，中间最细
        符号 & 含义 & 单位 \\% &表示分隔，\\表示换行
        \midrule[1pt]
        \addlinespace %\addlinespace 的作用是在表格行与行之间增加额外的垂直空白间距。
        \addlinespace
        \addlinespace
        \addlinespace
        \bottomrule[1.5pt]
    \end{tabularx}
\end{table}
\begin{enumerate}[label=(\arabic*), leftmargin=2.5em, itemsep=0.3em]
\item \textbf{只列重要符号。} 将正文中多次使用、含义关键的变量、参数和指标集中列出，临时变量一般不必放入符号说明。

\item \textbf{符号与含义一一对应。} 每个符号应说明其具体含义，必要时标明单位，避免只写模糊解释。

\item \textbf{全文保持统一。} 同一符号在不同章节中应表示同一含义，不要中途更换符号或一符多义。

\item \textbf{注意符号规范。} 变量通常使用斜体字母，向量、矩阵和集合可根据全文习惯统一表示。

\item \textbf{首次出现仍要解释。} 即使已经在符号说明表中列出，正文第一次使用时也建议简要说明。

\item \textbf{控制篇幅。} 符号说明应尽量控制在1页以内，避免把所有临时变量都列入表中。
\end{enumerate}
\textbf{符号说明的作用，是帮助读者快速查找变量含义，而不是代替正文解释。}

\textbf{写作原则：重要、统一、清楚、简洁，尽量不超过1页。}
      % 四、符号说明
\section{问题一模型的建立与求解}

\subsection{数据预处理}
\subsubsection{缺失值处理}
\subsubsection{数据清洗与标准化}
\begin{enumerate}[label=(\arabic*), leftmargin=2.5em, itemsep=0.3em]
\item \textbf{说明数据来源。} 明确数据由赛题附件、公开数据库、文献资料或自主采集获得；使用外部数据时，应在正文和参考文献中注明来源。

\item \textbf{先检查数据质量。} 建模前应分析数据是否存在缺失值、异常值、重复值、格式错误及量纲不一致等问题。

\item \textbf{处理方法要有依据。} 对缺失值填补、异常值剔除、数据平滑等操作，应说明采用的方法和理由，不能随意修改原始数据。

\item \textbf{根据模型需要进行转换。} 若不同指标的量纲、数量级或正负方向不一致，可进行正向化、标准化或归一化处理。

\item \textbf{保留有效信息。} 数据预处理应减少噪声和无效信息，但不能过度处理，更不能为了得到理想结果而人为调整数据。

\item \textbf{保证过程可复现。} 关键处理步骤、参数设置和处理前后的数据变化应交代清楚；较长的程序和完整数据可放入附录或支撑材料。
\end{enumerate}

\textbf{数据预处理的目的，是提高数据质量并满足模型要求，而不是改变数据所反映的客观规律。}

\textbf{常见处理内容}

\ding{172} \textbf{缺失值处理} 根据缺失比例和数据特点，采用删除、均值填补、插值或模型预测等方法。

\ding{173} \textbf{异常值处理} 可通过箱线图、(3$\sigma$) 原则、四分位距或专业阈值识别异常值，并结合实际背景判断是否修正或剔除。

\ding{174} \textbf{数据清洗与转换} 统一日期、单位和数据格式，删除重复记录，必要时进行对数变换、平滑处理或变量编码。

\ding{175} \textbf{无量纲化处理} 评价类和机器学习类问题中，可采用极差标准化、Z-score标准化等方法消除量纲影响。

\textbf{写作提示}

\begin{enumerate}[label=(\arabic*), leftmargin=2.5em, itemsep=0.3em]
\item 只写与后续建模有关的处理步骤，不必罗列所有软件操作。

\item 预处理方法应与数据特点相匹配，避免机械套用固定方法。

\item 建议通过表格或图形展示处理前后的主要变化，但图表不宜过多。

\item 若原始数据质量较好，也应简要说明已完成完整性和异常性检查。
\end{enumerate}

\textbf{写作原则：来源清楚、处理合理、过程透明、结果可信。}

\subsubsection{问题一求解思路}
% 原文档此处为 OLE 公式对象（Equation Editor 3.0），无法自动转为 LaTeX 数学源码。
% 已导出其预览图：gq_template_images/equation_ole.wmf（可用 Inkscape 等转 PNG 后替换下方 includegraphics）。
% \includegraphics[width=0.5\textwidth]{gq_template_images/equation_ole.wmf}
\begin{equation}                                % 行间带编号
  a^2 + b^2 = c^2
\end{equation}
\subsubsection{问题一模型建立}
\subsubsection{问题一模型求解与分析}
\begin{figure}[!h]
    \centering
    
\includegraphics[width=.6\textwidth]{数模加油站logo}
    \caption{数模加油站}
    \label{fig:circuit-diagram}
\end{figure}
\begin{enumerate}[label=(\arabic*), leftmargin=2.5em, itemsep=0.3em]
\item \textbf{模型应紧扣问题。} 根据问题分析和模型假设，明确决策变量、目标函数、约束条件或变量之间的数学关系。

\item \textbf{说明模型选择依据。} 简要解释所选模型为什么适合本题，避免只给出模型名称或机械套用现成方法。

\item \textbf{推导过程清晰。} 重要公式应说明来源和含义，公式之间要有逻辑联系，不能只堆砌公式而缺少文字解释。

\item \textbf{求解方法明确。} 说明采用的算法、软件及主要步骤，必要时给出参数设置、终止条件和计算流程。

\item \textbf{结果完整可信。} 根据题目要求给出数值、方案、图表或评价结果，并结合实际背景解释其含义。

\item \textbf{按问题分别展开。} 建议采用“问题一模型建立与求解”“问题二模型建立与求解”等小标题，使论文结构与赛题任务保持一致。
\end{enumerate}
\textbf{模型建立回答“数学关系是什么”，模型求解回答“如何得到结果”。}

\textbf{推荐写作结构}

\ding{172} \textbf{确定变量与参数} 说明输入数据、决策变量、状态变量及主要参数的含义。

\ding{173} \textbf{建立数学模型} 给出目标函数、约束条件、评价指标或变量关系，并说明公式的实际意义。

\ding{174} \textbf{选择求解方法} 介绍所用算法及选择理由，说明主要计算步骤和关键参数。

\ding{175} \textbf{展示求解结果} 使用表格或图形呈现核心结果，并围绕题目要求进行解释。

\textbf{写作提示}
\begin{enumerate}[label=(\arabic*), leftmargin=2.5em, itemsep=0.3em]
\item 公式应编号，正文引用时使用“式（1）”“式（2）”等规范表述。

\href{https://www.bilibili.com/video/BV1LpN2zQEo1?spm_id_from=333.788.videopod.sections&vd_source=b669ebd8251242404e23b3f8835d481e}{mathtype的安装与使用教程}

%\begin{enumerate}[label=(\arabic*), leftmargin=2.5em, itemsep=0.3em]
\item 成熟模型可以引用文献，但必须结合题目进行调整，不能直接照搬。

\item 程序代码不宜大段放入正文，核心流程可用流程图或伪代码表示，完整代码放入支撑材料。

\item 图表应具有标题、编号、坐标名称和单位，正文中应对重要图表进行分析。

\item 结果保留合理的小数位和有效数字，避免出现无意义的过高精度。

\item 不要只给计算结果，应说明结果反映的实际规律及其是否合理。
\end{enumerate}

\textbf{写作原则：模型合理、推导清楚、求解可复现、结果有解释。}
      % 五、问题一模型的建立与求解
\section{问题二模型的建立与求解}
\subsection{问题二求解思路}
\subsection{问题二模型建立}
\subsection{问题二模型求解与分析}
      % 六、问题二模型的建立与求解
\section{问题三模型的建立与求解}
\subsection{问题三求解思路}
\subsection{问题三模型建立}
\subsection{问题三模型求解与分析}
      % 七、问题三模型的建立与求解
\section{灵敏度分析}

\textbf{灵敏度分析}是通过改变模型中的关键参数或输入数据，观察模型输出及最终结论的变化程度，从而判断模型对参数扰动是否敏感、结果是否稳定；\textbf{误差分析}是比较模型计算结果与真实值、观测值或参考值之间的差异，并分析误差大小及其来源，从而评价模型的准确性。二者的主要区别在于：灵敏度分析关注“\textbf{参数变化会使结果改变多少}”，主要检验模型的稳定性和鲁棒性；误差分析关注“\textbf{模型结果与实际值相差多少}”，主要检验模型的准确性和可信度。

\subsection{误差分析}
\begin{enumerate}[label=(\arabic*), leftmargin=2.5em, itemsep=0.3em]
\item \textbf{检验模型可靠性。} 模型求解后，应分析参数变化、数据误差和模型假设对结果的影响，判断结论是否稳定可信。

\item \textbf{选择合适的检验方法。} 并非所有模型都必须同时进行多种检验，应根据模型类型和数据特点选择灵敏度分析、误差分析或其他检验方法。

\item \textbf{围绕关键结果展开。} 重点检验重要参数、核心指标和最终方案，不必对所有变量逐一分析。

\item \textbf{结果应有解释。} 不能只给出检验图表或误差数值，还应说明模型是否稳定、误差是否可接受以及主要影响因素是什么。
\end{enumerate}
\hspace*{\parindent} \textbf{模型检验的目的，是说明结果是否稳定、准确、可信，而不是简单增加图表和计算。}
%\hspace*{\parindent} 默认系统缩进

\textbf{一、灵敏度分析}
\begin{enumerate}[label=(\arabic*), leftmargin=2.5em, itemsep=0.3em]
\item \textbf{选择关键参数。} 优先分析取值具有不确定性、对模型结构影响较大或来源于人为设定的参数。

\item \textbf{合理设置变化范围。} 可在基准值附近按一定比例调整参数，如分别变化 (±5\%)、(±10\%) 或其他符合实际的范围。

\item \textbf{保持其他条件不变。} 每次改变一个参数，观察目标值、评价结果、排序或最优方案的变化；必要时可分析多个参数的共同影响。

\item \textbf{展示变化规律。} 可使用表格、折线图或柱状图表示参数变化与模型结果之间的关系。

\item \textbf{判断模型稳定性。} 若参数小幅变化时结果变化较小，说明模型稳定；若结果变化明显，应指出敏感参数并分析原因。
\end{enumerate}
\hspace*{\parindent} \textbf{灵敏度分析重点回答：参数发生变化时，模型结论是否仍然成立。}

\textbf{二、误差分析}
\begin{enumerate}[label=(\arabic*), leftmargin=2.5em, itemsep=0.3em]
\item \textbf{明确误差来源。} 误差可能来自原始数据、测量过程、缺失值处理、参数估计、模型简化或数值计算。

\item \textbf{选择误差指标。} 可根据问题类型采用绝对误差、相对误差、均方误差、平均绝对误差、拟合优度等指标。

\item \textbf{进行结果对比。} 将模型结果与真实值、历史数据、题目给定结果或其他合理方法的结果进行比较。

\item \textbf{分析误差大小。} 说明误差是否处于可接受范围，并结合实际背景解释偏差产生的原因。

\item \textbf{指出模型局限。} 对误差较大的情况，应说明模型在哪些条件下可能失效，以及可以如何改进。
\end{enumerate}
\hspace*{\parindent} \textbf{误差分析重点回答：模型结果与实际情况相差多少，误差从哪里产生。}

\textbf{写作提示}
\begin{enumerate}[label=(\arabic*), leftmargin=2.5em, itemsep=0.3em]
\item 灵敏度分析和误差分析应与具体模型结合，避免只介绍方法而没有实际计算。

\item 图表应突出关键变化和主要结论，不必展示大量重复结果。

\item 不能仅凭一次计算就说明模型稳定，应设置合理的参数变化范围或对照数据。

\item 若缺少真实数据，无法直接计算误差，应如实说明，并采用交叉验证、历史回代或情景分析等方式间接检验。
\end{enumerate}
\hspace*{\parindent} \textbf{写作原则：指标合理、过程清楚、结论明确、评价客观。}
   % 八、灵敏度分析
\section{模型评价与推广}
\begin{enumerate}
\item \textbf{客观评价模型。} 从模型假设、计算效率、稳定性、适用范围等方面分析模型的优点与不足，做到客观真实。

\item \textbf{结合模型特点。} 所有评价都应结合本文建立的模型展开，避免使用“精度高”“效果好”等空泛描述。

\item \textbf{提出改进方向。} 对存在的问题可简要说明今后的优化思路，但不必重新建立新模型。

\item \textbf{分析推广价值。} 说明模型是否能够应用于类似问题，以及推广过程中需要调整的内容。
\end{enumerate}
\hspace*{\parindent}\textbf{本部分主要体现对模型的全面认识，不是论文重点，做到分析客观、内容精炼即可。}


\textbf{一、模型优点}

\begin{enumerate}
\item 分析模型是否具有较好的\textbf{科学性和合理性}，能够较真实地描述实际问题。

\item 说明模型是否具有\textbf{较好的可解释性}，变量和参数具有明确的实际意义。

\item 评价模型的\textbf{计算效率}，分析算法是否收敛较快、计算复杂度是否较低。

\item 说明模型是否具有\textbf{较好的稳定性}，当参数或数据发生小范围变化时，结果是否保持稳定。

\item 分析模型是否具有\textbf{较好的通用性}，是否能够适用于同类问题。
\end{enumerate}

\textbf{二、模型缺点}

\begin{enumerate}
\item 分析模型假设是否过于理想化，对实际问题进行了哪些简化。

\item 说明模型对数据质量、参数选取或初始条件是否较为敏感。

\item 分析模型适用范围是否有限，在哪些条件下可能不再适用。

\item 指出模型是否忽略了部分影响因素，或某些复杂关系未能充分考虑。

\item 结合实际情况说明模型仍存在的不足，并提出后续改进方向。
\end{enumerate}

\textbf{三、模型推广}

\begin{enumerate}
\item 说明模型的核心思想是否可以应用于同类型问题。

\item 分析模型推广到其他地区、行业或研究对象时，需要调整哪些参数或约束条件。

\item 若更换数据来源或增加新的影响因素，应说明模型需要进行哪些修改。

\item 可以结合人工智能、大数据、优化算法等方法，简要说明模型未来的改进和应用方向。
\end{enumerate}
\paragraph{写作提示}
\begin{enumerate}
\item 优点和缺点应一一对应，评价要有依据，不宜过分夸大或刻意贬低。

\item 推广应立足于模型本身，不宜脱离论文内容进行过多延伸。

\item 不必每一点都详细展开，可结合实际模型选择其中2～4点进行说明。

\item 建议控制篇幅，一般\textbf{0.5～1页}即可。
\end{enumerate}
\paragraph{写作原则：评价客观、优缺点对应、推广合理、内容简洁。}
\subsection{模型的优点}
\begin{enumerate}
\item 模型充分结合实际，简化xx，xx，xx条件，考虑了诸多重要因素得到合理的模型，如:xx，xx，xx。这样得到的模型贴合实际，具有较高的应用价值，可以推广到xx；【模型的假设好】
\item 模型运用xx和xx思想，抓住影响xx问题的重要因素，将复杂的xx问题转化为简单的xx问题，合理设置参数，模型的输出结果符合题目要求，能解决实际问题；【模型的参数好】
\item 本文使用的xx算法具有xx，xx，xx等优点，对于求解xx模型非常适用；【模型的求解算法好】
\item 本文得到的xx(安排方案、策略)具有效率高、输出稳定、xx均衡等特点，基本不存在xx，xx，xx等问题，在现有条件下能有效提高生产效率。【模型的结果好】
\end{enumerate}
\subsection{模型的不足}
\begin{enumerate}
\item 实际应用中，xx和xx可能也是重要的因素，但本文未能考虑到这些因素的影响，一定程度上影响了模型的准确性；【模型不够好】
\item 本文提出的模型对于现有条件使用效果较好，由于时间问题没有对其他情况进行检验。对于其他情形(如:xx，xx)，可能无法达到较好的效果；【模型适用范围较窄】
\item 实际上xx，xx的影响不一定是线性的，而本文将其作为线性因子处理，忽略了边际效应的影响。【模型的因素有些不好】
\end{enumerate}
\subsection{模型的推广}
【推广：在xx方面，可以将xx参数替换成xx参数，从而解决xx问题；
改进：结合参考文献xx，进一步考虑xx的影响，从而得到更合理的模型；
这部分不用太多，至少4行】
    % 九、模型评价与推广
\section{AI使用声明}

【下面根据实际情况二者选一】

本参赛队在竞赛过程中未使用任何AI工具。

本参赛队在竞赛过程中使用了AI工具，主要用于【简要用途，如语言润色、代码调试等】，详细使用情况见支撑材料。
 %十、AI声明
\section{参考文献}
引用别人的成果或其他公开的资料(包括网上查到的资料)必须按照规定的参考文献的表述方式在正文引用处和参考文献中均明确列出。正文引用处用方括号标示参考文献的编号，如[1][3]等；引用书籍还必须指出页码。参考文献按正文中的引用次序列出，其中

书籍的表述方式为：

[编号] 作者，书名，出版地：出版社，出版年.

参考文献中期刊杂志论文的表述方式为：

[编号] 作者，论文名，杂志名，卷期号：起止页码，出版年.

参考文献中网上资源的表述方式为：

[编号] 作者，资源标题，网址，访问时间（年月日）.
\begin{thebibliography}{99}
\bibitem{} 汤国生.基于SWOT分析法的大学生数学建模创新实践基地建设探索——以江苏科技大学为例[J].高教学刊,2023,9(11):53-56.DOI:10.19980/j.CN23-1593/G4.2023.11.013.
\bibitem{} 谢雨婧,韩惠丽.北师大版初中数学教材中数学建模的多维度分析[J].教学与管理,2023(09):73-76.
\bibitem{} 刘志梅.数学建模与高职数学教学的深度融合[J].佳木斯职业学院学报,2023,39(03):152-154.
\bibitem{} 许亚桃,吴立宝.基于Delphi-AHP高中数学建模教学评价指标体系的研究[J].内江师范学院学报,2023,38(02):113-119.DOI:10.13603/j.cnki.51-1621/z.2023.02.018.
\bibitem{} 杨本朝,石雅男,段乾恒,李光松,于刚.大学生数学建模竞赛开展全周期教学实践探究[J].大学教育,2023(04):44-46.
\bibitem{} 黄健,徐斌艳.国际视野下数学建模教与学研究的发展趋势——基于第14届国际数学教育大会的分析[J].数学教育学报,2023,32(01):93-98.
\begin{figure}[!h]
    \centering
    \includegraphics[width=.6\textwidth]{示例图1}
    \caption{示例图1}
    \label{fig:circuit-diagram}
\end{figure}

\begin{figure}[!h]
    \centering
    \includegraphics[width=.6\textwidth]{示例图2}
    \caption{示例图2}
    \label{fig:circuit-diagram}
\end{figure}
\end{thebibliography} %十一、参考文献
%% ------------------------------------------------------------
%% 参考文献：GB/T 7714-2015 数字顺序制（gbt7714 宏包 + bibtex）
%% 只允许引用真实存在的文献；条目写入 references.bib
%% ------------------------------------------------------------
%%\nocite{*} % example.tex 原未使用 \cite，此处列出全部参考文献
%%\bibliographystyle{gbt7714-numerical}
%%\bibliography{references}

%% ------------------------------------------------------------
%% 附录：核心代码（listings 排版，见 mmastyle 的 mmacode 样式）
%% + AI 使用声明（国赛 risk_extra 硬性要求，确定性模板，勿删）
%% ------------------------------------------------------------
\appendix
\section{附录A 文件列表}
\begin{table}[H] % H强制原地，不再漂移
\centering %\centering —— 对齐命令（让表格居中）
\begin{tabular}{|c|c|} %{c c}则是三线表，没有竖框
\hline %在当前行的下方画一条贯穿整个表格的横线。
文件名 & 功能描述 \\ \hline
 &  \\ \hline
 &  \\ \hline
 &  \\ \hline
 &  \\ \hline
 &  \\ \hline
 &  \\ \hline
 &  \\ \hline
 &  \\ \hline
 &  \\ \hline
 &  \\ \hline
\end{tabular}
\end{table}


\begin{center}
\Large\textbf{附录B代码}
\end{center}

\subsection*{基于 xxxx 的第一问求解代码}
\begin{lstlisting}[language=Python]
# 你的Python示例代码
import numpy as np

def solve_question1():
    res = np.array([1,2,3])
    return res

if __name__ == "__main__":
    ans = solve_question1()
    print(ans)
\end{lstlisting}




%\begin{center}
\Large\textbf{代码}
\end{center}

\subsection*{基于 xxxx 的第一问求解代码}
\begin{lstlisting}[language=Python]
# 你的Python示例代码
import numpy as np

def solve_question1():
    res = np.array([1,2,3])
    return res

if __name__ == "__main__":
    ans = solve_question1()
    print(ans)
\end{lstlisting}
% \input{chapter/C_ai_declaration}

\end{document}
